\section{Seeding the textual dimensions}
\label{sec:seeding}

Appendix~\ref{app:seeds} gives the seed text and the behavioural readout for each
profile. The five attribute dimensions are properties of the people behind the agents, and a
harness cannot instantiate them. It can assert them---which is exactly what a
deployed agent card, reputation field, or relationship summary does. We therefore
seed each as context and measure whether behavior changes.


\paragraph{The origin condition is the cleanest instrument.} Four seeds change what is
stated about incentives or consequences. \dimname{Relationship origin} changes only how
the encounter is narrated: same directory, same grants, same knowledge, same store. If
the primary metrics move under that manipulation alone, relational framing does real
work; if they do not, richer metadata on a card is unlikely to do better, because it is
the same channel carrying more words.

\paragraph{Validity boundary.} \emph{Seeded trust is not trust.} We measure whether an
agent acts differently when told it is in a trusting, aligned, accountable or
long-standing relationship---the sensitivity of a runtime to asserted relational
context, which is what an agent network actually delivers to a model. We do not measure
what trust does in a human network, and no claim here should be read that way. If seeds
move behavior the paper reports elasticities usable in card and protocol design; if not,
it reports that relational metadata is inert without architectural backing.
